% Generated draft: requires mathematical and source review.
\documentclass[11pt]{article}
\usepackage[margin=25mm]{geometry}
\usepackage{fontspec}
\setmainfont{Latin Modern Roman}
\usepackage{amsmath,amssymb,mathtools,unicode-math}
\setmathfont{Latin Modern Math}
\usepackage{enumitem,needspace,xurl,hyperref}
\setlength{\parindent}{0pt}
\setlength{\parskip}{5pt}
\setlength{\emergencystretch}{3em}
\newcommand{\N}{\mathbb{N}}
\newcommand{\Z}{\mathbb{Z}}
\newcommand{\Q}{\mathbb{Q}}
\newcommand{\R}{\mathbb{R}}
\newcommand{\C}{\mathbb{C}}
\newcommand{\F}{\mathbb{F}}
\newcommand{\A}{\mathbb{A}}
\newcommand{\PP}{\mathbb{P}}
\newcommand{\E}{\mathbb{E}}
\newcommand{\eps}{\varepsilon}
\newcommand{\dd}{\,\mathrm{d}}
\begin{document}
\section*{UnsolvedMath ID: 30003852}
\textbf{Source problem code:} catalogue entry
\subsection{Definitions and mathematical statement}
The supplied registry record combines two independent questions and marks its statement as \texttt{unrecoverable}; it does not promote a verified clean statement. The formulations below preserve those two questions provisionally, with the unresolved terminology stated explicitly.

A group is finitely presented if it admits a presentation with finitely many generators and finitely many defining relations. A group \(G\) has type \(F_\infty\), written \(F^{\infty}\) in the source, if there exists an Eilenberg--Mac Lane CW complex \(K(G,1)\) whose \(n\)-skeleton is finite for every integer \(n\geq 0\). An embedding of groups means an injective group homomorphism.

For a homeomorphism \(g\) of an interval \(I\), define its moved set by
\[
M(g)=\{x\in I:g(x)\neq x\}.
\]
The usual one-bump condition is that \(M(g)\) is a single nonempty interval; agreement with the source's precise convention requires review. The source writes \(\mathrm{PLo}(I)\), customarily suggesting orientation-preserving piecewise-linear homeomorphisms of an interval, but does not specify \(I\), its boundary conditions, or the precise piecewise-linearity convention. Its title refers instead to piecewise-projective groups. These choices cannot be resolved from the supplied evidence. Write \(\operatorname{Fast}(S)\) solely as a placeholder for the source's undefined condition that \(S\) is a fast set; no mathematical definition of this predicate is supplied here.

\textbf{Question 1, conditional on recovering these conventions.} Is it true that every finitely presented subgroup \(G\leq\mathrm{PLo}(I)\) which admits a generating set \(S\subseteq G\) satisfying \(\operatorname{Fast}(S)\), with every member of \(S\) a one-bump element, has type \(F_\infty\)? Equivalently, the intended universal assertion appears to be
\[
\forall G\leq\mathrm{PLo}(I),\quad
\left[
\begin{gathered}
G\text{ is finitely presented},\\
\exists S\subseteq G:\ G=\langle S\rangle,\quad
\operatorname{Fast}(S),\\
\forall s\in S,\ M(s)\text{ is a single nonempty interval}
\end{gathered}
\right]
\Longrightarrow G\text{ has type }F_\infty.
\]
This is an interpretation of the supplied wording, not a verified recovery of the numbered source question. The record does not explicitly require the fast generating set itself to be finite, so that restriction has not been inserted.

\textbf{Question 2.} Let
\[
\Gamma=\left\langle a,b,c,d\ \middle|\ b^a=b^2,\ c^b=c^2,\ d^c=d^2,\ a^d=a^2\right\rangle
\]
be the Higman group appearing in the record, using the conjugation convention \(x^y=y^{-1}xy\). Does there exist a Lodha--Moore group \(L\) and an injective homomorphism \(\Gamma\hookrightarrow L\)? The source also asks the weaker target question: does an injective homomorphism
\[
\Gamma\hookrightarrow H(\R)
\]
exist, where \(H(\R)\) denotes Monod's group named in the source? The record supplies neither a definition of the allowed Lodha--Moore groups nor a definition of \(H(\R)\); their exact meanings and the claimed weaker-target relationship require source verification.
\subsection{Short English statement}
The record merges two questions: whether every finitely presented subgroup of the stated interval homeomorphism group that has fast one-bump generators is of type \(F_\infty\), and whether the displayed Higman group embeds into a Lodha--Moore group, or at least into Monod's \(H(\R)\). Essential conventions and target-group definitions are missing. These are provisional statement-only formulations requiring human mathematical review.
\subsection{Sources}
\begin{itemize}
\item Local UnsolvedMath registry: supplied complete record with permanent numeric ID \(30003852\), exact source problem code \texttt{catalogue entry}, and title \emph{Finitely Presented Piecewise-Projective Groups}. The record lists proposed year 2018 and supplies no named proposer. Its statement-verification field reports that the extraction merges two independent numbered questions, leaves the original record visible for provenance, and promotes no separate clean statement; the statement status is \texttt{unrecoverable}.
\item Supplied report citation: \emph{Cohomological and Metric Properties of Groups of Homeomorphisms of R} (2019), Oberwolfach Reports, pp. 1579--1633, DOI 10.4171/owr/2018/26. \url{https://doi.org/10.4171/owr/2018/26}. The supplied background describes Oberwolfach Reports as documenting mathematical workshops and their open-problem discussions. No exact page or separate question numbers for these two questions are supplied.
\item Supplied literature source: \emph{Oberwolfach Report 26/2018}. \url{https://ems.press/content/serial-article-files/46748?nt=1}.
\item Supplied literature source: Lodha, \emph{A nonamenable type \(F_\infty\) group}. \url{https://ems.press/journals/ggd/articles/13662}.
\item Literature assessment dated exactly \texttt{2026-08-21}: the supplied curation labels the questions open and states that known type \(F_\infty\) results for particular Lodha--Moore groups settle neither question. This is supplied dated curation derived from the OWR material, not an independently verified literature assessment. The background expressly requests independent verification of time-sensitive status and cited claims. No source was visited for this draft.
\end{itemize}
\end{document}
